\section{Position of the results and remaining mathematical boundaries}\label{sec:comparison}
\subsection{Standard tools and the new construction}
Conditional projection, nearest-center quantization, small-ball converses, and bounded-loss comparison are standard tools, not new foundations in themselves. The quantization background is developed in Graf and Luschgy~\cite{graf}; comparison of experiments in Blackwell~\cite{blackwell}, Le Cam~\cite{lecam}, and Torgersen~\cite{torgersen}. Predictive-state representations and future-equivalence constructions have substantial antecedents in Littman, Sutton, and Singh~\cite{psr}, Shalizi and Crutchfield~\cite{shalizi}, and kernel sufficiency treatments such as Fritz~\cite{fritz}. Proposition~\ref{prop:quotient} records the additional pointwise instrument and measurable-image hypotheses actually used here; it is not a claim to originate predictive sufficiency.

The mathematical contribution of Theorem~\ref{thm:inward} is the constructive implication from a static \emph{inward} cover and true-law raw moment bounds to a legal finite recursive state, with a converse stated in the geometry of the actual experiment. Inwardness controls the approximate trajectory's envelope, while the two-point inequality controls its separation from the true trajectory. Neither an existing recursive quantizer nor intermediate-to-terminal occupation domination is a premise. The proofs use standard $L^2$ and Lyapunov arguments; the claim concerns their specified causal-geometric implication and its verified realizations, not the discovery of these analytic inequalities.

Average contraction of random Lipschitz maps is classical; see Diaconis and Freedman~\cite{diaconis}. Our square-moment condition is stronger than a mere negative logarithmic average and is not a replacement for that theory. It is chosen because it controls accumulated squared task error under report laws that depend on the true predictive state. The additional envelope is needed when a finite code on an unbounded stable coordinate has resolution depending on location. The product compander supplies an explicit construction in that setting. It does not claim that every average-stable random system admits the same rate.

\subsection{Filtering approximation and causal coding}
Quantization methods for nonlinear filtering are developed by Pag\`es and Pham~\cite{pages}, Sellami~\cite{sellami}, and Feuer and Goodwin~\cite{feuer}. The present paper does not claim novelty for recursively quantizing a filter or for propagating approximation error. Its finite-state realization adds a matching converse for the \emph{actually acquired} posterior law, uniform over a declared time-varying class with intermittent mixing, and combines that converse with output-alphabet and calibration restrictions in one scored experiment. The lower uses only the genuine gate-one submeasure and a raw Jacobian calculation, not a global posterior-density hypothesis.

McDonald and Y\"uksel~\cite{mcdonald} establish true-law expected filter contraction through transition and observation Dobrushin coefficients. This is an important antecedent for using the physical report law rather than the incorrect filter's imagined law. Our finite-state calculation uses an elementary common-mass bound in log coordinates and additionally constructs an inward finite-label recursion.

Kara and Y\"uksel~\cite{kara}, in particular their Theorems~12, 16, and 17, relate finite-window control approximation to filter stability under specified discount and regularity conditions. Their control objective is broader than our fixed-exploration prediction objective. Our label budget is the cardinality of a recursive retained state, rather than a prescribed observation-window length; our continuous-report example adds an actual geometric lower bound. Neither statement implies the other without an additional conversion between legal policy classes and resources. Saldi, Y\"uksel, and Linder~\cite{saldi} likewise study near-optimal control through finite belief-model approximations. We do not reinterpret a same-controller comparison as an optimal-control theorem.

Nonanticipative rate-distortion theory, including Charalambous, Stavrou, and Ahmed~\cite{nonanticipative}, imposes causal realizability on reconstruction kernels. Its information objective is not identical to a hard bound on every retained-state alphabet, even though information inequalities can supply converse tools. Proposition~\ref{prop:interface} identifies a precise continuation requirement forcing disjoint label supports. It does not turn every simulator implementation into an information-theoretically necessary product overhead.

\subsection{Singular geometry and complementary results}
The geometry of Moran quantization has an extensive theory; Zhu~\cite{zhu-moran} studies asymptotic uniformity for such measures. The novelty asserted in Theorem~\ref{thm:singular} is not a new Cantor-dimension formula. Its object is a finite, physically acquired cylinder law transported recursively through expanding and non-erasing transitions. The common-oracle identity keeps its atomic acquisition variance, while the static completed-domain cover and raw pair drift construct the online machine without an invariant preparation comparison.

The companion mathematical supplement~\cite{supplement} preserves the earlier allocation-weighted causal resolvent, terminal-bridge, finite-chart, robust-mixture, single-probe, and soft-state information results. These remain complementary: for example, nonuniform component allocations or separately certified continuation-information budgets need not satisfy a uniform inward cover at a single common scale. Conversely the present transient raw-drift argument does not require their occupation-comparison certificate. We claim neither implication without checking the respective hypotheses.

\subsection{Scope}
The general theorem and the two class verifications are complete under their displayed assumptions. Several larger objectives remain distinct proof obligations. A necessary-and-sufficient characterization of when actual checkpoint geometry is recursively realizable is not proved. Nor is a joint optimal region for persistent state, temporary workspace, program description, computation time, unknown-kernel learning, and causal deficiency. The continuation cut and output-alphabet theorem give genuine converses for their specified restrictions, not that entire region. Adaptive exploration and error-dependent stopping can change both the acquired law and the conditional inequalities; they require new analysis rather than a formal minimization of the present fixed-protocol bound.

These are boundaries of the conclusions, not substitutions for the experimental problem. The results establish a general constructive route and two different raw-kernel verifications, while keeping the broader acquired-geometry and resource-transfer question mathematically explicit.
